% AdvancedOpticsQuality20261023
\section{波动光学}

\subsection{标量波、复振幅与傍轴近似}

标量波满足
\[
 \left(\frac1{c^2}\frac{\partial^2}{\partial t^2}-\nabla^2\right)
 A(\vect r,t)=0.
\]
在直角、柱和球坐标中，典型外行波分别为
\[
\begin{aligned}
 A_{\rm p}&=A_0\cos(\omega t-\vect k\cdot\vect r+\varphi),\\
 A_{\rm c}&=\frac{A_0}{\sqrt{\rho}}\cos(\omega t-k\rho+\varphi),\\
 A_{\rm s}&=\frac{A_0}{r}\cos(\omega t-kr+\varphi),
\end{aligned}
\qquad \frac{\omega}{|k|}=c.
\]
上式柱面波是 $k\rho\gg1$ 时的远场渐近式，严格的单频径向解为 Hankel 函数 $H_0^{(1)}(k\rho)$；球面式在离开源点处是严格外行解。柱面波和球面波的幅度衰减分别保证单位长度圆柱面、球面上的平均能流守恒。

\begin{center}\begin{minipage}{\linewidth}\centering
\centering
\begin{tikzpicture}[>=Latex,scale=.72]
 \begin{scope}[xshift=-4.7cm]
  \foreach \y in {-.9,-.3,.3,.9}{\draw[DeepBlue,thick] (-1.45,\y)--(1.45,\y);}
  \foreach \x in {-1.0,0,1.0}{\draw[->,Gold!85!black] (\x,-1.15)--(\x,1.45);}
  \node at (0,-1.55){平面波};
 \end{scope}
 \begin{scope}
  \foreach \r in {.35,.72,1.1,1.48}{\draw[DeepBlue,thick] (25:\r) arc[start angle=25,end angle=155,radius=\r];}
  \foreach \a in {35,65,95,125,150}{\draw[->,Gold!85!black] (0,0)--(\a:1.75);}
  \draw[RuleGray] (-.22,-.2)--(.22,.2);
  \node at (0,-1.55){柱面波截面};
 \end{scope}
 \begin{scope}[xshift=4.7cm]
  \foreach \r in {.38,.76,1.14,1.52}{\draw[DeepBlue,thick] (0,0) circle (\r);}
  \foreach \a in {20,80,145,210,275,335}{\draw[->,Gold!85!black] (0,0)--(\a:1.82);}
  \fill (0,0) circle(1.2pt);
  \node at (0,-1.95){球面波截面};
 \end{scope}
\end{tikzpicture}
\bookcaption{平面波、柱面波和球面波的波前与传播方向}
\label{advA5:fig:three-waves}
\end{minipage}\end{center}

统一采用光学符号约定
\[
 A(\vect r,t)=A(\vect r)\e^{-\ii\omega t},\qquad
 A(\vect r)=|A|\e^{\ii\phi(\vect r)},\qquad
 \vect k=\nabla\phi.
\]
单色场的空间部分满足 Helmholtz 方程
\[
 \boxed{(\nabla^2+k^2)A(\vect r)=0},\qquad k=\omega/c.
\]
在均匀、各向同性、线性无源介质中，电磁场各分量满足相应标量 Helmholtz 方程，势函数还须选用相容规范；非均匀介质不能无条件逐分量套用同一标量方程，
但偏振问题仍需保留矢量方向。

取主传播轴为 $z$，傍轴平面波写成
\[
 A(x,y,z)=A_0\e^{\ii kz}
 \e^{\ii k(\theta_xx+\theta_yy)-\ii kz(\theta_x^2+\theta_y^2)/2},\qquad
 |\theta_x|,|\theta_y|\ll1.
\]
物平面 $z=-L$ 上 $(X,Y)$ 点发出的球面波传播到像平面 $(x,y,0)$ 时，
\[
\begin{aligned}
 R&=\sqrt{L^2+(x-X)^2+(y-Y)^2}\\
 &\simeq L+\frac{(x-X)^2+(y-Y)^2}{2L},
\end{aligned}
\]
所以
\[
 A(x,y)\propto
 \e^{\ii k(x^2+y^2)/(2L)}
 \e^{-\ii k(Xx+Yy)/L}.
\]
图 \ref{advA5:fig:spherical-plane-approx} 中虚线是局部切平面；
图中 $X,Y,x,y$ 的位置与公式保持同一符号约定。

\begin{center}\begin{minipage}{\linewidth}\centering
\centering
\begin{tikzpicture}[>=Latex,scale=.76]
 \draw[->] (-5,0)--(5.4,0) node[right]{$z$};
 \draw[thick] (-4,-1.6)--(-4,1.8) node[above]{物平面};
 \draw[thick] (3.8,-1.6)--(3.8,1.8) node[above]{像平面};
 \coordinate (S) at (-4,.8);
 \coordinate (P) at (3.8,.42);
 \fill (S) circle(1.3pt) node[left]{$(X,Y)$};
 \fill (P) circle(1.3pt) node[right]{$(x,y)$};
 \draw[DeepBlue,thick,->] (S)--(P) node[midway,above]{$R$};
 \draw[DeepBlue!65] (S) ++(-65:1.5) arc[start angle=-65,end angle=65,radius=1.5];
 \draw[DeepBlue!65] (S) ++(-42:3.0) arc[start angle=-42,end angle=42,radius=3.0];
 \draw[DeepBlue!55,dashed] (2.82,-1.15)--(3.05,1.55)
   node[above left,text=TextGray,xshift=-6pt]{局部切波前};
 \draw[<->] (-4,-1.95)--(3.8,-1.95) node[midway,below]{$L$};
 \draw[Gold!80!black] (-3.15,.8) arc[start angle=0,end angle=-3,radius=.85];
 \node[Gold!70!black] at (-3.0,.57) {$\theta$};
\end{tikzpicture}
\bookcaption{球面波在像平面附近的傍轴展开与局部平面波近似}
\label{advA5:fig:spherical-plane-approx}
\end{minipage}\end{center}

当观测孔径半径为 $\rho$ 时，若
\[
 \rho^2\ll\lambda L,
\]
二次相位在孔径内可以忽略，得到远场平面波近似。若不满足此式，
保留二次相位即为 Fresnel 传播；这不是“误差项”，而是近场衍射和散焦的主要来源。

\subsection{相干与多光束干涉}

两复振幅叠加为 $A=A_1+A_2$，探测器测得时间平均光强
\[
 I=I_1+I_2+2\sqrt{I_1I_2}\,
 \operatorname{Re}\gamma_{12},
\qquad
 \gamma_{12}=
 \frac{\langle A_1^*A_2\rangle}
 {\sqrt{\langle|A_1|^2\rangle\langle|A_2|^2\rangle}}.
\]
$|\gamma|=1$ 为完全相干，$|\gamma|=0$ 为完全不相干。扩展光源限制空间相干性，
典型相干角量级为 $\Delta\theta\sim\lambda/b$；有限谱宽限制时间相干性，
相干长度量级为 $L_c\sim\lambda^2/\Delta\lambda$。

\begin{advexamplebox}{两界面的 Stokes 关系与等效反射、透射}\label{advA5:ex:multibeam}
三层无吸收介质有两个平行界面。正向通过第 $i$ 个界面的振幅反射、透射系数为
$r_i,t_i$，反向为 $r_i',t_i'$；层内单程相位为 $\delta$。

\begin{enumerate}
\item 从能流守恒和时间反演说明 Stokes 关系的正确复数形式；
\item 求整个双界面的等效振幅反射系数与透射系数。
\end{enumerate}

各介质中的通道均为传播波，选定同一偏振并按单位能流归一化，排除全反射的倏逝通道情形。

\QAOFrame{\QAOInterfaces{0}}{两平行界面的传播通道及单程相位；幅度系数使用同一能流规范。}{quality:A5:q7}
\end{advexamplebox}
\begin{advsolution}
若先把行波振幅按单位能流归一化，无吸收界面的散射矩阵
\[
 S_i=\begin{pmatrix}r_i&t_i'\\ t_i&r_i'\end{pmatrix}
\]
满足 $S_i^\dagger S_i=I$。互易、时间反演对称时还可取 $S_i=S_i^{\mathsf T}$。
因此一般的复数 Stokes 关系是
\[
 |r_i|^2+|t_i|^2=1,\qquad
 r_i^*t_i'+t_i^*r_i'=0.
\]
只有在系数可同时取实数并采用特定相位规范时，才化成
$r_i+r_i'=0$、$t_it_i'+r_i^2=1$。这修正了直接把“实系数关系”用于任意复系数的错误。

从第一界面反射后的各束构成几何级数：
\[
 r=r_1+t_1t_1'r_2\e^{2\ii\delta}
 \sum_{m=0}^{\infty}(r_1'r_2\e^{2\ii\delta})^m.
\]
故
\[
 r=r_1+\frac{t_1t_1'r_2\e^{2\ii\delta}}
 {1-r_1'r_2\e^{2\ii\delta}},
\qquad
 t=\frac{t_1t_2\e^{\ii\delta}}
 {1-r_1'r_2\e^{2\ii\delta}}.
\]
若采用实系数相位规范 $r_1'=-r_1$、$t_1t_1'=1-r_1^2$，
便化为
\[
 \boxed{r=\frac{r_1+r_2\e^{2\ii\delta}}
 {1+r_1r_2\e^{2\ii\delta}}},
 \qquad
 \boxed{t=\frac{t_1t_2\e^{\ii\delta}}
 {1+r_1r_2\e^{2\ii\delta}}}.
\]
$r_2=0$ 时退化为单界面 $r=r_1$；相位平均后交叉项消失，回到非相干多次反射。

\QAOFrame{\QAOInterfaces{1}}{两条透射路径与第一次内部反射；每一新增往返都贡献相同的复数乘子。}{quality:A5:a7}

图~\ref{quality:A5:a7}中第一束在界面1直接反射，第二束先进入中间层，再经界面2反射并从界面1退出。因此第二束的系数同时含 $t_1,t'_1,r_2$ 和往返相位；后续往返再乘 $r'_1r_2\e^{2\ii\delta}$。能流归一化后透射率才可直接写为 $|t|^2$。
\end{advsolution}

\begin{advexamplebox}{周期双层介质的允许带与禁带}\label{advA5:ex:photonic-crystal}
折射率 $n_1,n_2$、厚度 $d_1,d_2$ 的两层介质周期交错。
分析哪些频率的正入射光可以传播，哪些频率落入禁带。

介质均匀、无吸收、非磁且折射率为正，先讨论无限周期结构的 Bloch 模式。

\QAOFrame{\QAOStackQuestion}{无限周期双层结构；每一周期的两个折射率与厚度必须分别保留。}{quality:A5:q8}
\end{advexamplebox}
\begin{advsolution}
一周期 $\Lambda=d_1+d_2$ 的特征矩阵为 $M=M_2M_1$。对非磁介质正入射，
令相位厚度 $\delta_j=k_0n_jd_j$，光学导纳 $\eta_j=n_j$，单层矩阵
\[
 M_j=
 \begin{pmatrix}
 \cos\delta_j&\ii\sin\delta_j/\eta_j\\
 \ii\eta_j\sin\delta_j&\cos\delta_j
 \end{pmatrix}.
\]
Bloch 条件要求本征值为 $\e^{\pm\ii K\Lambda}$，且 $\det M=1$，故
\[
 \boxed{\cos K\Lambda=
 \cos\delta_1\cos\delta_2
 -\frac12\left(\frac{\eta_1}{\eta_2}
 +\frac{\eta_2}{\eta_1}\right)
 \sin\delta_1\sin\delta_2}.
\]
当右端绝对值不超过 $1$ 时 $K$ 为实数，是允许带；超过 $1$ 时
$K$ 有非零虚部，场随周期数指数衰减，是禁带。四分之一波堆
$n_1d_1=n_2d_2=\lambda_0/4$ 在中心波长处右端
$=-(n_1/n_2+n_2/n_1)/2<-1$（只要 $n_1\ne n_2$），所以形成反射禁带。
$n_1=n_2$ 时右端化为 $\cos(\delta_1+\delta_2)$，禁带消失，极限正确。

\QAOFrame{\QAOStackAnswer}{定量示例：$n_1/n_2=2$、$n_1d_1=n_2d_2=\lambda_0/4$；阴影内矩阵迹超出传播区间。}{quality:A5:a8}

图~\ref{quality:A5:a8}的阴影由 $|\operatorname{tr}M/2|>1$ 确定，对应 Bloch 特征值一增一减。有限厚度的无吸收叠层仍可有隧穿透射，禁带不是吸收带；带边的重根也应单独看待，不能从 $K$ 实数就断言任意有限器件都完全透射。
\end{advsolution}

\subsection{衍射、路径积分与光栅}

\begin{advtheorem}{Kirchhoff 衍射积分}\label{advA5:thm:kirchhoff}
设 $A$ 在体积 $V$ 内满足 Helmholtz 方程，Green 函数
$G(\vect r,\vect r')=\e^{\ii kR}/R$，$R=|\vect r-\vect r'|$。
若边界由孔径面和远处闭合面组成，则
\[
 A(\vect r)=\frac1{4\pi}\iint_{\partial V}
 \left(G\frac{\partial A}{\partial n'}
 -A\frac{\partial G}{\partial n'}\right)\dd S'.
\]
在 Kirchhoff 边界近似下，孔径 $S$ 上可写成
\[
 \boxed{A(\vect r)\simeq\frac1{\ii\lambda}
 \iint_S A(\vect r')\frac{\e^{\ii kR}}{R}
 \frac{\cos\theta+\cos\theta'}2\,\dd S'}.
\]
\end{advtheorem}

\noindent\textit{证明。}
对 $A$ 与 $G$ 应用 Green 第二恒等式：
\[
 \iiint_V(A\nabla'^2G-G\nabla'^2A)\dd V'
 =\iint_{\partial V}\left(
 A\frac{\partial G}{\partial n'}
 -G\frac{\partial A}{\partial n'}\right)\dd S'.
\]
$G$ 满足 $(\nabla'^2+k^2)G=-4\pi\delta(\vect r-\vect r')$，
$A$ 满足齐次 Helmholtz 方程，代入后左端给出 $-4\pi A(\vect r)$，
移项即得第一式。远处半球在 Sommerfeld 辐射条件下无贡献。
在孔径上把 $A$ 和法向导数取为无屏时的值，在遮挡部分取零；
对 $G$ 的法向导数保留主导的 $\ii k\cos\theta/R$ 项，
整理得到倾斜因子 $(\cos\theta+\cos\theta')/2$。这也说明 Kirchhoff 边界条件是近似，
不能同时作为严格的 Dirichlet 与 Neumann 条件。

在透镜后焦面观察且忽略缓变倾斜因子时，得到 Fraunhofer 公式
\[
 \boxed{A(\theta_x,\theta_y)=
 \frac{\e^{\ii kf}}{\ii\lambda f}
 \iint A(x',y')\e^{-\ii k(\theta_xx'+\theta_yy')}
 \dd x'\dd y'}.
\]
因此远场衍射花样是孔径复振幅的二维 Fourier 变换。

\begin{advtheorem}{光传播的路径积分表述}\label{advA5:thm:path-integral}
在高频标量近似下，传播核可形式地写为
\[
 K(B,A)=\frac1{\mathcal N}\int_{\gamma:A\to B}
 \exp\!\left[\ii k_0\int_\gamma n(\vect r)\,\dd s\right]
 \mathcal D\gamma .
\]
短波极限 $k_0\to\infty$ 中，主要贡献来自光程驻值路径，即费马光线。
\end{advtheorem}

\noindent\textit{证明思路。}
把传播距离分成许多薄层，在每一层使用 Fresnel 核
\[
 K(\vect r_{j+1},\vect r_j)\propto
 \exp\!\left[\ii k_0n_j
 \left(\Delta z+\frac{|\Delta\vect r_\perp|^2}{2\Delta z}\right)\right],
\]
并对所有中间横向位置积分。令层数趋于无穷，就得到对路径的泛函积分。
当 $k_0$ 很大时，非驻相邻路径的相位快速转动而相消；驻相条件
$\delta\int n\,\dd s=0$ 正是费马原理。对驻值路径作二次展开，
高斯积分给出几何扩散振幅和 Maslov 相位，说明“只剩一条光线”只是领头阶近似。

\begin{advexamplebox}{2020级联双缝的中心光强与谱选择}\label{advA5:ex:cascade-slits}
单缝 $P_0$ 后依次放置 $P_1,\ldots,P_{2020}$ 共 $2020$ 个等间距双缝面，
相邻面间距 $L$，每个双缝间距 $D$，末端屏中心为 $O$。
所有缝均等宽，双缝关于主轴对称，满足傍轴条件。$P_0$ 到 $P_1$ 距离为 $L_0$，最后双缝到屏距离为 $L_1$。若把这些双缝全换成同宽、位于轴上的单缝，$O$ 点光强为 $I_0$。比较不同波长时取相同的参考强度 $I_0$，即比较相对增强因子 $I_O/I_0$。本题只比较理想传播模型的相对光强；实际多次遮挡会使绝对透射功率很小。

\begin{enumerate}
\item 求双缝级联时 $O$ 点光强；
\item 取 $L_0=L=L_1=5\,\mathrm{cm}$、$D=0.4\,\mathrm{mm}$，
求可见光中使中心光强最大的波长；
\item 求中心光强不小于最大值一半的波长范围；
\item 说明白光下条纹的颜色、间距和级次。
\end{enumerate}

\QAOFrame{\begin{tikzpicture}[>=Latex,font=\small]
 \draw[thick](-5,-1.1)--(-5,-.06) (-5,.06)--(-5,1.1)node[above]{$P_0$};
 \foreach \x/\t in {-3/P_1,-1/P_2,2/P_{2020}}{
  \draw[thick](\x,-1.1)--(\x,-.45) (\x,-.31)--(\x,.31) (\x,.45)--(\x,1.1)node[above]{$\t$};
 }
 \draw[thick](4,-1.1)--(4,1.1)node[above]{$P$};
 \fill(4,0)circle(1.4pt)node[right]{$O$};
 \draw[DeepBlue!60](-5,0)--(-3,.38)--(-1,-.38);
 \draw[DeepBlue!60](-5,0)--(-3,-.38)--(-1,.38);
 \draw[DeepBlue!60](2,-.38)--(4,0) (2,.38)--(4,0);
 \node at(.5,0){$\cdots$};
 \foreach \a/\b/\lab in {-5/-3/L_0,-3/-1/L,2/4/L_1}{
  \draw[<->](\a,-1.45)--(\b,-1.45)node[midway,below]{$\lab$};}
 \draw[<->](-3.5,-.38)--(-3.5,.38)node[midway,left]{$D$};
\end{tikzpicture}}{单缝入射、双缝级联和末端中心；省略号仅省略中间面，不改变面数。\label{advA5:fig:cascade-source}}{quality:A5:q9}
\end{advexamplebox}
\begin{advsolution}
采用窄缝离散传播模型：每缝内的振幅和相位近似均匀，忽略斜程的振幅差及单缝包络变化。若缝宽为 $a$，除 $a\ll D$ 外，还需 $2\pi aD/(\lambda L)\ll1$。
在相邻双缝面之间，同侧传播距离为 $L$，交叉传播距离在傍轴下为
$L+D^2/(2L)$。略去公共相位，单步传播矩阵为
\[
 T=\begin{pmatrix}1&\e^{\ii\delta}\\
 \e^{\ii\delta}&1\end{pmatrix},\qquad
 \delta=\frac{2\pi}{\lambda}\frac{D^2}{2L}.
\]
入射与中心读出都对应对称本征向量 $(1,1)^{\mathsf T}$，
其本征值为 $1+\e^{\ii\delta}
=2\e^{\ii\delta/2}\cos(\delta/2)$。
与全改单缝的振幅比较，第一处分束和末端合束再给出一个因子 $2$，
故
\[
 \frac{A_O}{A_{O,0}}
 =2^{2020}\e^{\ii\Phi}
 \cos^{2019}\!\frac{\delta}{2},
\]
\[
 \boxed{\frac{I_O}{I_0}
 =4^{2020}
 \cos^{4038}\!\left(\frac{\pi D^2}{2\lambda L}\right)}.
\]

数值上
\[
 \Delta\ell=\frac{D^2}{2L}=1.6000\,\mu\mathrm m.
\]
最大值要求 $\Delta\ell/\lambda=m$ 为整数。按可见范围约
$380$--$780\,\mathrm{nm}$，有
\[
 \boxed{\lambda=533.33\,\mathrm{nm}\ (m=3)},\qquad
 \boxed{\lambda=400.00\,\mathrm{nm}\ (m=4)}.
\]
若采用不同的可见光边界约定，边界附近的 $400\,\mathrm{nm}$ 是否计入应随约定说明。

半高条件为
\[
 \left|\cos\frac{\pi\Delta\ell}{\lambda}\right|
 \ge 2^{-1/4038}.
\]
令
$q=\arccos(2^{-1/4038})=0.018528\ldots$，每个整数 $m$ 周围满足
\[
 \left|\frac{\pi\Delta\ell}{\lambda}-m\pi\right|\le q.
\]
于是两段可见范围近似为
\[
 \boxed{532.29\text{--}534.38\,\mathrm{nm}},\qquad
 \boxed{399.41\text{--}400.59\,\mathrm{nm}}.
\]
\textbf{谱选择与空间条纹应分开。}最后一组双缝的两出射振幅始终相等。屏上纵坐标为 $y$ 时，末段光程差为 $Dy/L_1$，所以忽略单缝衍射包络的变化，有
\[
 \frac{I(y,\lambda)}{I_0(\lambda)}=
 4^{2020}\cos^{4038}\!\left(\frac{\pi\Delta\ell}{\lambda}\right)
 \cos^2\!\left(\frac{\pi Dy}{\lambda L_1}\right).
\]
第一个因子选择窄波段，第二个因子给出普通双缝空间条纹：
\[
 y_k=k\frac{\lambda L_1}{D},\qquad
 \Delta y_{533}=66.667\,\mu\mathrm m,\quad
 \Delta y_{400}=50.000\,\mu\mathrm m.
\]
绿光的第 $3j$ 级与紫光的第 $4j$ 级在 $y=200j\,\mu\mathrm m$ 处重合。白光经级联后主要留下这两个窄波段；各自形成相应间距的亮暗纹，重合处颜色叠加，其余处随两分量比例改变颜色。中心的实际颜色、相对亮度还取决于入射谱和 $I_0(\lambda)$，不能只由几何量确定。有限谱宽会使远处条纹逐渐模糊。级联幂次使\emph{谱峰}变窄，并不把最后双缝的单色空间因子变成同样的高次幂。

\QAOFrame{\QAOCascadeAnswer}{左为级联的谱选择，右为最终双缝的单色空间条纹；两个不同变量上的峰宽不能混同。}{quality:A5:a9}

图~\ref{quality:A5:a9}左图以某一谱峰附近的相位失谐为横坐标，展示高次幂选谱；右图固定波长，只剩末端双缝的 $\cos^2$ 空间因子。两者的自变量分别包含 $\lambda$ 和 $y$，因此窄谱峰不能直接当作窄空间条纹。题图~\ref{quality:A5:q9}中的2020个双缝面只有2019次双缝到双缝传播。
\end{advsolution}

对周期为 $d$、缝宽 $a$、缝数 $N$ 的矩形光栅，
\[
 I(\theta)=I_1
 \left(\frac{\sin\alpha}{\alpha}\right)^2
 \left(\frac{\sin N\beta}{\sin\beta}\right)^2,
\quad
 \alpha=\frac{\pi a}{\lambda}\sin\theta,\quad
 \beta=\frac{\pi d}{\lambda}\sin\theta.
\]
主极大满足 $d\sin\theta_j=j\lambda$，强度随 $N^2$ 增长，
相邻零点给出的角宽约为
$2\lambda/(Nd\cos\theta_j)$。

\begin{advexamplebox}{用 Fresnel 双棱镜阵列模拟光栅}\label{advA5:ex:biprism-grating}
实验室只有狭缝、对称双凸薄透镜（两面曲率半径 $R=30\,\mathrm{cm}$、折射率 $n_0=1.5$）和许多薄 Fresnel 双棱镜（单侧劈角 $\alpha=30'$、折射率 $n=1.005$）。拟在 $\lambda=500\,\mathrm{nm}$ 时获得一级角位置 $\theta_1=30^\circ$，且后焦面主极大两侧第一零点间宽度不大于 $0.1\,\mathrm{mm}$。按原题要求讨论至少需要多少个双棱镜以及各仪器的放置方法，并检查所需近似和器材条件。
\QAOFrame{\begin{tikzpicture}[scale=.8]
 \draw[thick](-4,-1.2)--(-4,1.2)node[above]{狭缝 $S$};
 \foreach \x in {-1.8,-1.4,-1}{\filldraw[fill=DeepBlue!10](\x,-1)--(\x+.2,0)--(\x,1)--(\x-.1,0)--cycle;}
 \node at(-1.4,1.5){双棱镜组};
 \filldraw[thick,fill=DeepBlue!8](1,-1.3)..controls(.7,-.65)and(.7,.65)..(1,1.3)..controls(1.3,.65)and(1.3,-.65)..cycle;
 \node at(1,1.65){透镜 $L$};
 \draw[thick](4,-1.3)--(4,1.3)node[above]{屏};
 \draw[guide](-4.3,0)--(4.3,0);
 \node at(0,-1.8){器件顺序示意；棱脊横移量与间距尚须确定};
\end{tikzpicture}}{狭缝、横移双棱镜组、透镜与屏的器件顺序；单元数仍待求。}{quality:A5:q10}
\end{advexamplebox}
\begin{advsolution}
\textbf{先区分题目所给条件与理想阵列模型。}
透镜焦距 $f=R/[2(n_0-1)]=0.300\,\mathrm m$；单侧小劈角的偏折量为
$\delta=(n-1)\alpha=4.363\times10^{-5}\,\mathrm{rad}$。
题给 $30^\circ$ 与全程傍轴的高精度使用并不相容。以下先给出原题形式近轴模型中的构造，再说明它不能证明原器材条件下的全局最小数目。

若等效透射场由 $N$ 个等宽、等幅、等相位的平移单元组成，周期为 $d$，则近轴阵列因子为
\[
 \left|\sum_{j=0}^{N-1}e^{-ikjd\theta}\right|^2
 =\frac{\sin^2(Nkd\theta/2)}{\sin^2(kd\theta/2)}.
\]
一级主峰位置与两侧第一零点间的宽度给出
\[
 d=\frac{\lambda}{\theta_1}=0.9549\,\mu\mathrm m,
 \qquad w=\frac{2f\lambda}{Nd}=\frac{2f\theta_1}{N},
 \qquad \boxed{N\ge3142}.
\]
这里 $\theta_1$ 以弧度代入。$N$ 是场单元数，尚不是棱镜数。

\textbf{一个有明确附加条件的 $N-1$ 片方案。}
将 $N-1$ 片薄双棱镜等效贴合于同一横截面，棱脊横移至
\[
 b_j=(j-N/2)d\quad(j=1,\ldots,N-1),
 \qquad a=\frac{d}{2\delta}=10.94\,\mathrm{mm},
\]
其中 $a$ 为狭缝到该截面的距离。另须有公共孔径选出总宽 $Nd$ 的 $N$ 个等宽区段，并在这些区段上可近似等幅照明。薄相位模型给出
\[
 \Phi(y)=\frac{ky^2}{2a}-k\delta\sum_{j=1}^{N-1}|y-b_j|.
\]
第 $q=0,\ldots,N-1$ 个区段的中心为 $c_q=[q-(N-1)/2]d$。因该区间内绝对值之和的斜率为 $2q-(N-1)$，配方得
\[
 \Phi(y)=\frac{k(y-c_q)^2}{2a}+C,
\]
各段常数 $C$ 相同：两相邻段在公共边界上连续，且边界距两段中心均为 $d/2$。故各区段确为同一场分布的平移，透镜后焦面得到“单元包络乘阵列因子”。透镜放在阵列后，屏置于后焦面；其有效口径还须接收所需全部单元。
在这一限定构型内，每片双棱镜只增加一个相位折点，要产生 $N$ 个区段至少需 $N-1$ 片，因此形式近轴结果为
\[
 \boxed{N-1=3141\ \text{片（限定构型及公共孔径假设下）}}.
\]
同轴重叠而不横移的任意多片双棱镜仍仅给两段相位，不能按片数声称形成 $2^n$ 束。

\textbf{为什么这不是原题无条件的唯一数值答案。}
原器材清单没有规定公共孔径、可否贴合及横移、照明均匀性和大角度成像规律，因而上述方案有明确的附加条件。若保留精确的 $30^\circ$，并指定平面屏坐标 $Y=f\tan\theta$，则
\[
 d=\frac{\lambda}{\sin\theta_1}=1.000\,\mu\mathrm m,
 \quad w\simeq\frac{2f\lambda}{Nd\cos^3\theta_1},
 \quad N\ge4619.
\]
此数值仅是指定大角度屏映射下的理想阵列要求，不能继续直接套用前述近轴棱镜构造。原题若要有可实现且唯一的“最少片数”，必须补足这些条件；不能混用精确光栅角与 $Y=f\theta$ 给出混合近似的束数结论。

\QAOFrame{\QAOBiprismAnswer}{四个平移相位单元的简化示意；各段抛物线相同，公共边界上相位连续。}{quality:A5:a10}

图~\ref{quality:A5:a10}中的棱脊 $b_j$ 是相位斜率的折点，而每个区段的中心是抛物线的极小点。公共孔径必须切出这些完整区段，才能将透射场写成相同单元的平移求和；仅把棱镜同轴堆放不会产生该图的多区段结构。
\end{advsolution}

Fresnel 波带法把波前按到观察点的光程差每增加 $\lambda/2$ 分带。
在傍轴球面上，带状面元面积的增加恰好补偿传播距离变化，使相邻半波带振幅近似相等、相位相反。

\begin{center}\begin{minipage}{\linewidth}\centering
\centering
\begin{tikzpicture}[scale=.9]
 \begin{scope}
  \draw[thick,domain=-40:40,variable=\t]plot({-3+2*cos(\t)},{2*sin(\t)});
  \coordinate(P)at(2,0);\fill(P)circle(1.6pt)node[right]{$P$};
  \foreach \t in {-35,-20,-7,7,20,35}{\draw[DeepBlue!70]({-3+2*cos(\t)},{2*sin(\t)})--(P);}
  \node at(0,-1.8){同一球面上的面元至 $P$};
 \end{scope}
 \begin{scope}[xshift=6cm]
  \foreach \m in {1,...,6}{\draw[DeepBlue] (0,0)circle({.62*sqrt(\m)});}
  \draw[Gold!80!black,->](0,0)--({.62*sqrt(4)*cos(35)},{.62*sqrt(4)*sin(35)})node[above right]{$r_4$};
  \node at(0,-1.8){近轴半波带：$r_m\propto\sqrt m$};
 \end{scope}
\end{tikzpicture}
\bookcaption{Fresnel 波带与圆孔半波带的几何关系}
\label{advA5:fig:fresnel-zones}
\end{minipage}\end{center}

若把每个半波带的复振幅首尾相接，离散求和在复平面上来回折返；
对于直边衍射的一维二次相位积分，连续累积振幅形成 Cornu 螺旋；圆孔环带的面积权重不同，不把两种几何的积分曲线等同。图
\ref{advA5:fig:phasor-sum} 同时显示这两种合成方法。

\begin{center}\begin{minipage}{\linewidth}\centering
\centering
\begin{tikzpicture}[scale=1]
\begin{scope}\draw[axis](-.3,0)--(2.7,0)node[right]{$\operatorname{Re}A$};\draw[axis](0,-.3)--(0,2)node[above]{$\operatorname{Im}A$};
\draw[DeepBlue,->](0.000000,0.000000)--(1.163562,1.570739);
\draw[DeepBlue,->](1.163562,1.570739)--(1.074157,0.755514);
\draw[DeepBlue,->](1.074157,0.755514)--(1.114611,1.383667);
\draw[DeepBlue,->](1.114611,1.383667)--(1.090363,0.853532);
\draw[DeepBlue,->](1.090363,0.853532)--(1.106948,1.320797);
\draw[DeepBlue,->](1.106948,1.320797)--(1.094692,0.898263);
\draw[Gold!80!black,very thick,->](0,0)--(1.094692,0.898263)node[right]{$\sum A_m$};\node at(1,-.8){相邻半周期分组};\end{scope}
\begin{scope}[xshift=6cm]\draw[axis](-2,0)--(2,0);\draw[axis](0,-2)--(0,2);\draw[DeepBlue,thick]plot coordinates {(-1.332586,-1.091889) (-1.326154,-1.037393) (-1.307278,-0.985865) (-1.277074,-0.940047) (-1.237228,-0.902309) (-1.189888,-0.874542) (-1.137532,-0.858078) (-1.082825,-0.853647) (-1.028482,-0.861363) (-0.977129,-0.880748) (-0.931180,-0.910780) (-0.892734,-0.949963) (-0.863487,-0.996418) (-0.844671,-1.047991) (-0.837025,-1.102356) (-0.840777,-1.157129) (-0.855662,-1.209977) (-0.880957,-1.258709) (-0.915525,-1.301368) (-0.957890,-1.336299) (-1.006309,-1.362200) (-1.058851,-1.378157) (-1.113489,-1.383665) (-1.168173,-1.378622) (-1.220915,-1.363318) (-1.269857,-1.338401) (-1.313330,-1.304839) (-1.349907,-1.263867) (-1.378436,-1.216934) (-1.398069,-1.165636) (-1.408273,-1.111665) (-1.408830,-1.056738) (-1.399829,-1.002550) (-1.381645,-0.950715) (-1.354913,-0.902724) (-1.320495,-0.859907) (-1.279441,-0.823401) (-1.232949,-0.794131) (-1.182323,-0.772791) (-1.128929,-0.759843) (-1.074157,-0.755514) (-1.019381,-0.759807) (-0.965925,-0.772512) (-0.915030,-0.793223) (-0.867835,-0.821364) (-0.825347,-0.856210) (-0.788431,-0.896914) (-0.757797,-0.942535) (-0.733992,-0.992065) (-0.717404,-1.044457) (-0.708257,-1.098646) (-0.706621,-1.153579) (-0.712419,-1.208231) (-0.725442,-1.261626) (-0.745356,-1.312852) (-0.771724,-1.361076) (-0.804016,-1.405553) (-0.841628,-1.445633) (-0.883900,-1.480766) (-0.930125,-1.510506) (-0.979575,-1.534511) (-1.031503,-1.552537) (-1.085169,-1.564441) (-1.139840,-1.570169) (-1.194811,-1.569755) (-1.249405,-1.563312) (-1.302987,-1.551021) (-1.354968,-1.533129) (-1.404811,-1.509933) (-1.452031,-1.481778) (-1.496200,-1.449042) (-1.536948,-1.412134) (-1.573963,-1.371482) (-1.606989,-1.327526) (-1.635823,-1.280713) (-1.660318,-1.231489) (-1.680375,-1.180295) (-1.695939,-1.127561) (-1.707001,-1.073701) (-1.713589,-1.019113) (-1.715765,-0.964170) (-1.713625,-0.909226) (-1.707287,-0.854605) (-1.696895,-0.800609) (-1.682611,-0.747508) (-1.664612,-0.695548) (-1.643087,-0.644947) (-1.618234,-0.595893) (-1.590257,-0.548551) (-1.559362,-0.503059) (-1.525757,-0.459530) (-1.489647,-0.418054) (-1.451235,-0.378700) (-1.410718,-0.341516) (-1.368287,-0.306531) (-1.324126,-0.273757) (-1.278410,-0.243188) (-1.231304,-0.214807) (-1.182964,-0.188582) (-1.133537,-0.164468) (-1.083157,-0.142411) (-1.031950,-0.122350) (-0.980030,-0.104212) (-0.927501,-0.087920) (-0.874458,-0.073391) (-0.820984,-0.060534) (-0.767154,-0.049258) (-0.713034,-0.039465) (-0.658682,-0.031057) (-0.604147,-0.023932) (-0.549470,-0.017986) (-0.494687,-0.013115) (-0.439826,-0.009213) (-0.384911,-0.006173) (-0.329959,-0.003887) (-0.274983,-0.002250) (-0.219995,-0.001152) (-0.164999,-0.000486) (-0.110000,-0.000144) (-0.055000,-0.000018) (0.000000,0.000000) (0.055000,0.000018) (0.110000,0.000144) (0.164999,0.000486) (0.219995,0.001152) (0.274983,0.002250) (0.329959,0.003887) (0.384911,0.006173) (0.439826,0.009213) (0.494687,0.013115) (0.549470,0.017986) (0.604147,0.023932) (0.658682,0.031057) (0.713034,0.039465) (0.767154,0.049258) (0.820984,0.060534) (0.874458,0.073391) (0.927501,0.087920) (0.980030,0.104212) (1.031950,0.122350) (1.083157,0.142411) (1.133537,0.164468) (1.182964,0.188582) (1.231304,0.214807) (1.278410,0.243188) (1.324126,0.273757) (1.368287,0.306531) (1.410718,0.341516) (1.451235,0.378700) (1.489647,0.418054) (1.525757,0.459530) (1.559362,0.503059) (1.590257,0.548551) (1.618234,0.595893) (1.643087,0.644947) (1.664612,0.695548) (1.682611,0.747508) (1.696895,0.800609) (1.707287,0.854605) (1.713625,0.909226) (1.715765,0.964170) (1.713589,1.019113) (1.707001,1.073701) (1.695939,1.127561) (1.680375,1.180295) (1.660318,1.231489) (1.635823,1.280713) (1.606989,1.327526) (1.573963,1.371482) (1.536948,1.412134) (1.496200,1.449042) (1.452031,1.481778) (1.404811,1.509933) (1.354968,1.533129) (1.302987,1.551021) (1.249405,1.563312) (1.194811,1.569755) (1.139840,1.570169) (1.085169,1.564441) (1.031503,1.552537) (0.979575,1.534511) (0.930125,1.510506) (0.883900,1.480766) (0.841628,1.445633) (0.804016,1.405553) (0.771724,1.361076) (0.745356,1.312852) (0.725442,1.261626) (0.712419,1.208231) (0.706621,1.153579) (0.708257,1.098646) (0.717404,1.044457) (0.733992,0.992065) (0.757797,0.942535) (0.788431,0.896914) (0.825347,0.856210) (0.867835,0.821364) (0.915030,0.793223) (0.965925,0.772512) (1.019381,0.759807) (1.074157,0.755514) (1.128929,0.759843) (1.182323,0.772791) (1.232949,0.794131) (1.279441,0.823401) (1.320495,0.859907) (1.354913,0.902724) (1.381645,0.950715) (1.399829,1.002550) (1.408830,1.056738) (1.408273,1.111665) (1.398069,1.165636) (1.378436,1.216934) (1.349907,1.263867) (1.313330,1.304839) (1.269857,1.338401) (1.220915,1.363318) (1.168173,1.378622) (1.113489,1.383665) (1.058851,1.378157) (1.006309,1.362200) (0.957890,1.336299) (0.915525,1.301368) (0.880957,1.258709) (0.855662,1.209977) (0.840777,1.157129) (0.837025,1.102356) (0.844671,1.047991) (0.863487,0.996418) (0.892734,0.949963) (0.931180,0.910780) (0.977129,0.880748) (1.028482,0.861363) (1.082825,0.853647) (1.137532,0.858078) (1.189888,0.874542) (1.237228,0.902309) (1.277074,0.940047) (1.307278,0.985865) (1.326154,1.037393) (1.332586,1.091889)};
\draw[Gold!80!black,very thick,->](-1.573963,-1.371482)--(0.712419,1.208231)node[pos=.55,above left=4pt]{$A$};\fill(-1.573963,-1.371482)circle(1.5pt)(0.712419,1.208231)circle(1.5pt);\node at(0,-2.35){直边衍射的 Fresnel 积分};\end{scope}\end{tikzpicture}
\bookcaption{离散半波带合成与连续复振幅积分}
\label{advA5:fig:phasor-sum}
\end{minipage}\end{center}

\begin{advexamplebox}{阶梯折射率相位光栅的后焦面分布}\label{advA5:ex:gradient-grating}
光栅由 $2N$ 个宽度为 $a$、厚度为 $d$ 的透明条带组成。每一条带内部
$\dd n/\dd y=0$，相邻条带折射率增加 $\alpha a$。
波数为 $k$ 的平面波正入射，随后通过焦距 $F$ 的透镜。
求后焦面坐标 $Y$ 上的光强因子，并判断何时花样关于焦点对称。

\QAOFrame{\begin{tikzpicture}[scale=.9]
 \foreach \i in {0,...,7}{\pgfmathtruncatemacro{\shade}{8+5*\i}\filldraw[fill=DeepBlue!\shade,draw=RuleGray](0,{-1.2+.3*\i})rectangle(.5,{-1.2+.3*(\i+1)});}
 \draw[axis](0,-1.4)--(0,1.7)node[above]{$y$};
 \foreach \h in {-.9,0,.9}{\draw[DeepBlue,->](-1.6,\h)--(-.15,\h);}
 \filldraw[thick,fill=DeepBlue!8](2.5,-1.4)..controls(2.3,-0.7)and(2.3,0.7)..(2.5,1.4)..controls(2.7,0.7)and(2.7,-0.7)..cycle;
 \draw[thick](5,-1.5)--(5,1.5)node[above]{后焦面};
 \draw[guide,DeepBlue](.55,.9)--(2.5,.9)--(5,0);
 \draw[guide,DeepBlue](.55,-.9)--(2.5,-.9)--(5,0);
 \draw[<->](.75,-1.2)--(.75,-.9)node[right]{$a$};
 \draw[<->](0,-1.7)--(.5,-1.7)node[midway,below]{$d$};
 \node at(2.5,2.2){条带沿 $y$ 并列；虚线仅示透镜轴上聚焦};
\end{tikzpicture}}{折射率台阶沿横向分布；每条带内部均匀，透镜用于后焦面观察。}{quality:A5:q11}
\end{advexamplebox}
\begin{advsolution}
令后焦面的横向空间频率
$q=kY/F$。第 $m$ 条带（$m=-N,\ldots,N-1$）的附加相位为
$kd(n_0+m\alpha a)$。略去公共相位后，
\[
\begin{aligned}
 A(Y)&\propto
 \sum_{m=-N}^{N-1}
 \int_{ma}^{(m+1)a}
 \e^{\ii kd\alpha am}\e^{-\ii qy}\,\dd y\\
 &\propto
 \sinc\!\left(\frac{qa}{2}\right)
 \sum_{m=-N}^{N-1}\e^{\ii m\xi},
\end{aligned}
\]
其中
\[
 \xi=kd\alpha a-qa.
\]
有限几何级数给出
\[
 \boxed{I(Y)=I_c
 \sinc^2\!\left(\frac{kYa}{2F}\right)
 \frac{\sin^2(N\xi)}{\sin^2(\xi/2)}}.
\]
第一因子是单条带包络，第二因子是 $2N$ 条带结构因子。
当每条折射率台阶带来的相位增量为半周的整数倍
\[
 \boxed{kd\alpha a=\pi\ell,\qquad \ell\in\mathbb Z}
\]
时，利用结构因子 $G(\xi)=\sin^2(N\xi)/\sin^2(\xi/2)$ 的偶性和 $2\pi$ 周期性，可得 $G(\pi\ell-qa)=G(\pi\ell+qa)$，故整个花样关于焦点对称。奇数半周与整周的峰位不同，但都具有该对称性。一般相位台阶使主峰偏向
$Y_0=F d\alpha$（模倒易周期），即形成闪耀光栅。
令 $\alpha=0$ 可恢复 $2N$ 个相邻同相条带，即总宽 $2Na$ 的矩形孔径，量纲和极限均正确。

\QAOFrame{\QAOGratingAnswer}{取 $2N=8$：同相时向量相加，第一零点时向量闭合；每条带的公共包络另乘。}{quality:A5:a11}

图~\ref{quality:A5:a11}用八条带说明结构因子：$\xi=0$ 时所有相量同向，$\xi=\pi/4$ 时围成闭合八边形。实际强度还要乘单条带的 $\sinc^2$ 包络，不能把相量合成的主峰位置直接当作全部花样的唯一最大值。$\alpha=0$ 时各条带首尾相接且同相，实际退化为总宽 $2Na$ 的单个矩形孔径，而非缝间含不透光间隔的光栅。
对 $2N\ge2$，上述对称条件也是必要的。令 $M=2N$、$\varphi=kd\alpha a$，结构因子展开为 $G(\xi)=M+2\sum_{j=1}^{M-1}(M-j)\cos(j\xi)$。要求 $G(\varphi-u)=G(\varphi+u)$ 对所有 $u$ 成立，正弦级数的每个系数必须为零；其中 $j=1$ 的系数非零，便有 $\sin\varphi=0$，即 $\varphi=\pi\ell$。公共包络仅在离散点为零，不改变这一整体对称条件。
\end{advsolution}

\subsection{偏振密度矩阵}

% Completeness20261024: coherency matrix, its decomposition and an analyzer calculation.
在固定横截面上，用复振幅描述准单色光的两个横向分量，写成
$\vect A=(A_x,A_y)^{\mathsf T}$。把强度的比例常数吸收进振幅定义，便有
$I=\langle|A_x|^2+|A_y|^2\rangle$；尖括号表示探测时间内的平均。
稳定偏振光的两个分量有确定的振幅比和相位差，Jones向量即可描述；
若它们随时间随机变化，必须先平均振幅的乘积，而不能先平均振幅再平方。
相干矩阵定义为
\[
 P=\langle\vect A\vect A^\dagger\rangle
 =\begin{pmatrix}
 \langle|A_x|^2\rangle&\langle A_xA_y^*\rangle\\
 \langle A_yA_x^*\rangle&\langle|A_y|^2\rangle
 \end{pmatrix},\qquad I=\operatorname{tr}P.
\]
非对角元保留两个方向的相关性及相对相位。由定义直接得到 $P^\dagger=P$。
对任意复向量 $\vect z$，
\[
 \vect z^\dagger P\vect z
 =\left\langle|\vect z^\dagger\vect A|^2\right\rangle\ge0,
\]
所以 $P$ 半正定，两个本征值均非负。这一性质保证任意检偏方向测出的强度不会为负。

\paragraph{为什么能够分成完全偏振光和非偏振光}
设 $I>0$，将 $P$ 对角化，记本征值 $\lambda_1\ge\lambda_2\ge0$，
对应单位正交Jones向量 $\vect e_1,\vect e_2$。利用
$\vect e_1\vect e_1^\dagger+\vect e_2\vect e_2^\dagger=I_2$，有
\[
 \begin{aligned}
 P&=\lambda_1\vect e_1\vect e_1^\dagger
       +\lambda_2\vect e_2\vect e_2^\dagger\\
  &=(\lambda_1-\lambda_2)\vect e_1\vect e_1^\dagger+\lambda_2 I_2.
 \end{aligned}
\]
第一项沿固定偏振态 $\vect e_1$，强度为 $I_{\rm p}=\lambda_1-\lambda_2$；
第二项在任何正交偏振基中都相同，是强度为 $I_{\rm u}=2\lambda_2$ 的完全非偏振分量。
两项的强度相加等于 $I$。因此偏振度为
\[
 \mathcal P=\frac{I_{\rm p}}I
 =\frac{\lambda_1-\lambda_2}{\lambda_1+\lambda_2}
 =\sqrt{1-\frac{4\det P}{(\operatorname{tr}P)^2}}.
\]
纯偏振光的 $\lambda_2=0$，即 $\det P=0$、$\mathcal P=1$；
完全非偏振光则有 $\lambda_1=\lambda_2=I/2$，即 $P=(I/2)I_2$。
零光强时偏振度不定义，不能把零矩阵仅凭行列式为零称为纯偏振光。

\paragraph{从矩阵算出检偏器的读数}
图~\ref{comp:A5:analyzer}在横截面内标出检偏器透光轴
$\vect e_\theta=(\cos\theta,\sin\theta)^{\mathsf T}$。
理想线偏振片透过的振幅是 $A_\theta=\vect e_\theta^\dagger\vect A$，故
\[
 I_\theta=\langle|A_\theta|^2\rangle
 =\vect e_\theta^\dagger P\vect e_\theta
 =P_{xx}\cos^2\theta+P_{yy}\sin^2\theta
   +2\operatorname{Re}P_{xy}\sin\theta\cos\theta.
\]
\begin{center}\begin{minipage}{.94\linewidth}\centering
\begin{tikzpicture}[>=Latex,font=\small]
 \draw[->] (-1.5,0)--(1.8,0) node[right]{$x$};
 \draw[->] (0,-1.05)--(0,1.65) node[above]{$y$};
 \draw[very thick,DeepBlue,<->] (-.9,-.63)--(1.25,.875);
 \draw[->] (.65,0) arc[start angle=0,end angle=35,radius=.65];
 \node at (.9,.25) {$\theta$};
 \node[above right] at (1.05,.86) {$\vect e_\theta$};
 \node[below,align=center] at (0,-1.25) {传播方向垂直纸面\\透光轴与 $x$ 轴夹角为 $\theta$};
 \draw[->,thick] (2.35,.35)--(3.65,.35);
 \node[draw,rounded corners=1pt,align=center,text width=3.6cm] at (5,.35)
 {$I_\theta=\vect e_\theta^\dagger P\vect e_\theta$\\探测器记录平均强度};
\end{tikzpicture}
\figcap{偏振矩阵的检偏测量：图中的方向向量同时进入振幅投影和强度二次型。}
\label{comp:A5:analyzer}
\end{minipage}\end{center}

例如，把总强度 $I_0$ 的光分解为强度 $0.6I_0$ 的 $x$ 线偏振光和强度
$0.4I_0$ 的非偏振光，两部分互不相干。各自的相干矩阵相加，得到
\[
 P=0.6I_0\begin{pmatrix}1&0\\0&0\end{pmatrix}
   +0.2I_0 I_2
 =I_0\begin{pmatrix}0.8&0\\0&0.2\end{pmatrix}.
\]
由图~\ref{comp:A5:analyzer}的投影式，
\[
 \frac{I_\theta}{I_0}=0.8\cos^2\theta+0.2\sin^2\theta
 =0.5+0.3\cos2\theta.
\]
图~\ref{comp:A5:analyzer-result}中 $\theta=0$ 透过 $0.8I_0$，
$\theta=\pi/2$ 仍透过 $0.2I_0$：后者来自非偏振分量的一半。
本例本征偏振是线偏振，故旋转线检偏器的反差
$(I_{\max}-I_{\min})/(I_{\max}+I_{\min})=0.6$ 恰等于偏振度。
这个结论不能推广为任意光的测量公式：纯圆偏振光通过所有线检偏方向都透过一半，
却有 $\mathcal P=1$；一般椭圆偏振态还需波片才能测出虚部相关信息。
\begin{center}\begin{minipage}{.94\linewidth}\centering
\begin{tikzpicture}[>=Latex,font=\small,x=1.6cm,y=2.5cm]
 \draw[->] (0,0)--(3.5,0) node[right]{$\theta$};
 \draw[->] (0,0)--(0,1.02) node[above]{$I_\theta/I_0$};
 \draw[DeepBlue,thick,domain=0:3.141593,samples=100]
   plot(\x,{.5+.3*cos(2*\x*180/pi)});
 \draw[dashed,TextGray] (0,.2)--(3.141593,.2);
 \draw[dashed,TextGray] (0,.8)--(3.141593,.8);
 \node[left] at (0,.2) {$0.2$};\node[left] at (0,.8) {$0.8$};
 \node[below] at (0,0) {$0$};
 \node[below] at (1.570796,0) {$\pi/2$};\node[below] at (3.141593,0) {$\pi$};
 \fill[DeepBlue] (1.570796,.2) circle[radius=1.2pt];
\end{tikzpicture}
\figcap{上述混合光的检偏曲线：非偏振分量使最小透射强度保持为正。}
\label{comp:A5:analyzer-result}
\end{minipage}\end{center}

\paragraph{偏振元件怎样作用于相干矩阵}
若线性光学元件的Jones矩阵为 $J$，则 $\vect A'=J\vect A$，从定义即得
$P'=JPJ^\dagger$。理想无损波片的 $J$ 为酉矩阵，保留 $P$ 的本征值、总强度和偏振度，
但可以改变偏振椭圆。以快慢轴为坐标轴、忽略公共相位，可按相位约定写成
\[
 J_{\rm ret}=\begin{pmatrix}\e^{\ii\delta}&0\\0&1\end{pmatrix}.
\]
$\delta=\pi$ 为半波片，$\delta=\pi/2$ 为四分之一波片。
例如等幅同相输入 $(1,1)^{\mathsf T}/\sqrt2$ 经过后一种矩阵变为
$(\ii,1)^{\mathsf T}/\sqrt2$，两个分量相差四分之一周期，成为圆偏振光。
其旋向还须连同传播方向和时间因子约定说明。


\subsection{晶体光学}

各向异性线性介质满足
\[
 \vect D=\boldsymbol\varepsilon\vect E.
\]

\begin{advtheorem}{无耗散介质的介电张量对称性}\label{advA5:thm:epsilon-symmetric}
在无色散或准静态、无耗散的可逆线性介质中，若电场能密度是状态函数，则
\[
 \boxed{\boldsymbol\varepsilon^{\mathsf T}
 =\boldsymbol\varepsilon}.
\]
\end{advtheorem}

\noindent\textit{证明。}
从零场准静态建立电场，电场能密度的增量为
$\dd u=\vect E\cdot\dd\vect D$。若 $u$ 是状态函数，定义电共能
$u^*(\vect E)=\vect E\cdot\vect D-u$，则
\[
 \dd u^*=\vect D\cdot\dd\vect E,\qquad
 D_i=\frac{\partial u^*}{\partial E_i}.
\]
因此二阶偏导可交换给出
\[
 \varepsilon_{ij}=\frac{\partial D_i}{\partial E_j}
 =\frac{\partial^2u^*}{\partial E_j\partial E_i}
 =\frac{\partial^2u^*}{\partial E_i\partial E_j}
 =\varepsilon_{ji}.
\]
对这种线性介质才进一步有 $u=u^*=\tfrac12\vect E\cdot\vect D$。
若存在吸收，频域介电张量可为复数；互易性与无耗散是不同条件。
反 Hermitian 部分决定单色场的平均吸收，色散介质的储能还涉及介电响应的频率导数，不能简单照搬上述实对称、无色散结论。

实对称介电张量可正交对角化为
$\operatorname{diag}(\varepsilon_x,\varepsilon_y,\varepsilon_z)$。
三者相等为各向同性晶体；两者相等为单轴晶体；三者互异为双轴晶体。
相位法线方向 $\hat{\vect p}=\vect k/k$ 与能流方向
$\hat{\vect r}\parallel\vect S$ 一般不重合，所以相速度和射线速度需分别讨论。

\begin{advtheorem}{Fresnel 相速度面与射线速度面}\label{advA5:thm:fresnel-crystal}
在主轴坐标中令
$v_i=1/\sqrt{\mu\varepsilon_i}$。平面波的相速度 $v_p$ 与法线方向余弦
$p_i$ 满足
\[
 \boxed{
 \frac{p_x^2}{v_p^2-v_x^2}
 +\frac{p_y^2}{v_p^2-v_y^2}
 +\frac{p_z^2}{v_p^2-v_z^2}=0}.
\]
射线速度 $v_r$ 与能流方向余弦 $r_i$ 满足对偶式
\[
 \boxed{
 \frac{r_x^2}{v_r^{-2}-v_x^{-2}}
 +\frac{r_y^2}{v_r^{-2}-v_y^{-2}}
 +\frac{r_z^2}{v_r^{-2}-v_z^{-2}}=0}.
\]
\end{advtheorem}

上述分式须先乘去分母再按连续极限解释。相速度的完整多项式为 $\sum_i p_i^2\prod_{j\ne i}(v_p^2-v_j^2)=0$；沿主轴或出现相等本征值时，不能在分式形式中直接约去零因子而丢失偏振支。单轴情形可因式分解为 $[v_p^2-v_o^2][v_p^2-v_o^2\cos^2\theta-v_e^2\sin^2\theta]=0$，两支都须保留。射线速度的对偶分式亦须作相同处理。

\noindent\textit{证明。}
对 $\e^{\ii(\vect k\cdot\vect r-\omega t)}$ 代入 Maxwell 方程，
\[
 \vect k\times(\vect k\times\vect E)
 +\omega^2\mu\boldsymbol\varepsilon\vect E=0.
\]
写成分量并令 $v_p=\omega/k$：
\[
 \left(\varepsilon_i-\frac1{\mu v_p^2}\right)E_i
 +\frac{p_i}{\mu v_p^2}(\hat{\vect p}\cdot\vect E)=0.
\]
消去 $E_i$，要求非零解的行列式为零，整理即得第一式。
能流方向与波矢法线面互为 Legendre 对偶；用群速度
$\vect v_g=\nabla_{\vect k}\omega$ 求相速度面的法向，
再消去 $p_i$ 得第二式。

对单轴晶体，取 $v_x=v_y=v_o$、$v_z=v_e$。一支为寻常光
$v_{po}=v_o$，电场垂直主截面；另一支为非寻常光
\[
 \boxed{v_{pe}^2=v_o^2\cos^2\theta+v_e^2\sin^2\theta},
\]
其中 $\theta$ 为波法线 $\vect k$ 与光轴的夹角；其 $\vect k$ 与能流一般不平行。界面折射时，频率和切向波矢连续；
因此应先在相速度面上求允许 $\vect k$，再由曲面的法向求射线方向，
不能把各向同性 Snell 定律直接套给能流方向。
